\documentclass{article}
\usepackage{amsmath}
\usepackage{xcolor}
\begin{document}

\textbf{Cap 01 - Introduccion a las Ecuaciones Diferenciales}

\textbf{Definicion}.-

Una Expresion que contiene dentro de si diferenciales(Derivadas de una mas variables dependientes) como tal, cuyo objetivo es determinar una funcion solucion que satisfaga dicha expresion. 




\begin{gather*}
x+y+3z = 55 \\
\text{tratar de adivinar los posibles valores para} \\
x = ?   \rightarrow   2 \\
y = ?   \rightarrow  10 \\
z = ?   \rightarrow  5 \\
2+10+3(5) = 55 \\
x^{2}+3x-5=0   \rightarrow  x=? \\
\vspace{0.5em} \\
\text{Una E.D. es algo similar, donde en vez de tener variables reales} \\
\text{tenemos como incognita, }\text{\colorbox[RGB]{248,231,28}{$\text{derivadas}$}} \\
\frac{dy}{dx} + 5x + y = 0 \\
\text{que funcion de x, satisface la diferencial? de tal manera que se cumpla la expresion dada} \\
\vspace{0.5em} \\
y = x^{2}+3   \rightarrow  f(x) \implies  y(x)=? \\
y' = 2x \\
\vspace{0.5em} \\
 2x + 5x + y = 0  \rightarrow \text{ no satisface}
\end{gather*}


\textbf{Clasificacion es de las Ecuaciones Diferenciales}

Clasificacion por Tipo

a) Ecuaciones Diferenciales Ordinarias $\rightarrow$  \colorbox[RGB]{248,231,28}{no hay fraciones parciales}


\begin{gather*}
\frac{dy}{dx} + 5x + y = 0         \frac{d^{2}y}{dx^{2}} + \frac{dy}{dx}  - 5y = 0       \frac{d^{2}y}{dx^{2}} + \frac{dz}{dx}  - 5e^{x} = 3x+y 
\end{gather*}




b) Ecuaciones Diferenciales No Ordinarias$\rightarrow$  Estan presentes las derivadas Parciales


\begin{gather*}
\frac{\partial ^{2}u}{\partial x^{2}}+\frac{\partial ^{2}u}{\partial y^{2}} = 0       \frac{\partial ^{2}u}{\partial x^{2}}= \frac{\partial ^{2}u}{\partial t^{2}} -2
\end{gather*}


\textbf{Grado de Una Ecuacion Diferencial} $\rightarrow$  Asociado al \colorbox[RGB]{248,231,28}{Exponente Mayor} presente dentro de la E.D. que representa el orden de la E.D.


\begin{gather*}
(\frac{d^{2}y}{dx^{2}})^{3} + (\frac{dy}{dx})^{2}  - 5y = 0           \frac{d^{3}y}{dx^{3}} + (\frac{dy}{dx})^{3} +3x\text{\colorbox[RGB]{248,231,28}{$($}}\text{\colorbox[RGB]{248,231,28}{$\frac{d^{3}y}{dx^{3}}$}}\text{\colorbox[RGB]{248,231,28}{$)$}}\text{\colorbox[RGB]{248,231,28}{$^{2}$}}- 5y = 0  \\
\text{E.D. de Grado 3                           E.D. de Grado 2                       }
\end{gather*}


\textbf{Orden de una E.D.} $\rightarrow$  Representa la mayor derivada presente dentro de la E.D.


\begin{gather*}
\text{\colorbox[RGB]{248,231,28}{$\frac{d^{4}y}{dx^{4}}$}}+5y=x                     \frac{dy}{dx}+ \text{\colorbox[RGB]{248,231,28}{$\frac{d^{3}y}{dx^{3}}$}} - 5y = 0  \\
\text{E.D. de Orden 4                        E.D. de Orden 3}
\end{gather*}




\textbf{E.D. Lineales}

Apliquemos la analogia de un polinomio de grado n para entender el concepto


\begin{gather*}
x^{2}+3x+5 = 0 \\
Ax^{2}+Bx+C = 0  \rightarrow \text{ 2do Grado} \\
Ax^{4}+Bx^{3}+Cx^{2}+Dx+E = 0  \rightarrow \text{ 4to Grado} \\
x^{4}-12x^{3}+\frac{1}{2}x^{2}+\sqrt{168}x+8 = 0 \\
\vspace{0.5em} \\
\text{Todas estas expresiones pueden ser resolubles por Ruffini}(\text{ Se puede programar} \\
\text{un algoritmo que haga dicha tarea})
\end{gather*}


En ecuaciones diferenciales, existe un formato similar, su forma generica esta dada por:


\begin{gather*}
a_{n}(x)\frac{d^{n}y}{dx^{n}}+a_{n-1}(x)\frac{d^{n-1}y}{dx^{n-1}}+......a_{1}(x)\frac{dy}{dx}+ a_{0}(x)y = g(x) \\
\text{Ejemplo} \\
5x\frac{d^{4}y}{dx^{4}}+senx\frac{d^{3}y}{dx^{3}}+e^{3x}\frac{d^{2}y}{dx^{2}}+\frac{dy}{dx^{}} +\sqrt{5}y=x
\end{gather*}





\end{document}
